\paragraph{Schistosomiasis-related fibrosis and ultrasound grading.}
Schistosome infection and established organ morbidity require complementary assessment: parasitological, serological, and molecular tests address infection, while imaging characterizes hepatosplenic abnormalities \citep{weerakoon2015diagnosis,sah2015imaging}. Liver biopsy provides histological information but is invasive, motivating non-invasive assessment; the limitations of fibrosis tests also require care when interpreting their outputs \citep{bravo2001biopsy,patel2020limitations}. The target in this study is an expert-assigned ultrasound appearance score.

For \emph{S. japonicum}, the anatomical distribution of fibrosis is particularly relevant. A longitudinal cohort linked ultrasound-detectable parenchymal fibrosis to infection history and documented limited reversal of severe changes after treatment \citep{carlton2010morbidity}. The Basel protocol distinguishes portal and interseptal patterns when assessing Asian schistosomiasis \citep{richter2025basel}. Together with the imaging literature \citep{sah2015imaging}, these findings motivate preserving anatomical context and visible regional appearances when predicting severity.

Computational assessment has progressed from B-mode texture features and support vector machines \citep{yeh2003fibrosis} to multiparametric ultrasomics \citep{li2019ultrasomics}, schistosomiasis-focused radiomics \citep{guo2024radiomics}, and convolutional classification with METAVIR-derived targets \citep{lee2020ultrasound}. SFibAI supplies the fine-grained score representation used here \citep{xu2026sfibai}. Acquisition, labels, and cohorts differ across these studies. Our direct comparison therefore uses SFibAI on the same cohort splits, while investigating how view and lesion-location cues can participate in grading.

\paragraph{Anatomical context and regional evidence learning.}
Multitask learning shares representations across related objectives \citep{caruana1997multitask}; SonoNet demonstrates standard-plane recognition and anatomical localization in fetal ultrasound \citep{baumgartner2017sononet}. Regional aggregation offers a complementary way to connect local information to an image-level decision. Attention-based MIL learns aggregation weights \citep{ilse2018attention}; pathology methods extend this idea with instance clustering in CLAM \citep{lu2021clam}, instance relationships in TransMIL \citep{shao2021transmil}, and additive local contributions in Additive MIL \citep{javed2022additive}. FiLM provides a general mechanism for conditioning feature transformations \citep{perez2018film}. VCRE-Fib draws on shared learning, conditioning, and aggregation to interpret regional features in their acquisition-view context before combining them with global judgment. Its regions belong to one ultrasound image, and its supervised view and localization outputs also enter the grading computation.

\paragraph{Weak localization and inspectable predictions.}
BoxSup and BoxInst learn segmentation from boxes without pixel-level masks \citep{dai2015boxsup,tian2021boxinst}. Here, boxes mark local abnormal patches and may not enclose the entire lesion, so responses beyond them remain candidates for clinical review. Concept bottleneck models expose supervised intermediate concepts \citep{koh2020concept}; VCRE-Fib exposes view and localization outputs while retaining a global route. Grad-CAM, Grad-CAM++, and integrated gradients provide post-hoc attribution approaches \citep{selvaraju2017gradcam,chattopadhay2018gradcampp,sundararajan2017axiomatic}. Our trained localization participates in regional aggregation, but this role alone does not establish faithful attribution. Saliency sanity checks motivate evaluating explanations beyond visual plausibility \citep{adebayo2018sanity}. The present clinician-annotation comparisons support inspection of spatial correspondence.

